% !TEX root =./Admin/ECE346_Syllabus_Fall 2024.tex

\begin{\myfontsize}

%\renewcommand{\arraystretch}{1.25}
\newcolumntype{b}{X} % big column
%\newcolumntype{m}{>{\hsize=.2\hsize}X} % medium column, .2 width of big
%\newcolumntype{s}{>{\hsize=.1\hsize}X} % small column, .1 width of big
\newcounter{objnum}
\setcounter{objnum}{1}
\newcounter{chapnum}
\setcounter{chapnum}{1}

%-----------------------------------------------------------
\def\chrowcolor{green!10!white}

\begin{center}
\begin{xltabular}{0.9\textwidth}{|c|X |} %columns are sb
\hline
\textbf{Obj}   & \textbf{Description} \\ \hline\hline\endhead % & \textbf{Format}\\ \hline\hline\endhead

\rowcolor{\chrowcolor}\multicolumn{2}{|c|}{Chapter {\thechapnum} - Introduction}
\\ \hline

\thechapnum.\theobjnum
& I can articulate what I will learn and how I will be assessed in this course.
\stepcounter{objnum}
\\\hline

\thechapnum.\theobjnum
& I can articulate practical applications of electromagnetics.
\setcounter{objnum}{1}
\stepcounter{chapnum}
\\\hline

%%%%%%%%%%%%%%%%%%%%%%%%%%%%%%%%%%%%%%%%%%%%%%%%%%%%%%%%%%%%%%%

\rowcolor{\chrowcolor}\multicolumn{2}{|c|}{Chapter {\thechapnum} - Fields}
\\\hline

\thechapnum.\theobjnum
& I can recall the definition of scalar fields, vector fields, phasors, and Euler's identity.
\stepcounter{objnum}
\\\hline 

\thechapnum.\theobjnum
& I can explain the basic terms and quantities used to describe electromagnetic fields. 
\setcounter{objnum}{1}
\stepcounter{chapnum}
\\\hline 

%%%%%%%%%%%%%%%%%%%%%%%%%%%%%%%%%%%%%%%%%%%%%%%%%%%%%%%%%%%%%%%

\rowcolor{\chrowcolor}\multicolumn{2}{|c|}{Chapter {\thechapnum} - Transmission Lines}
\\\hline

\thechapnum.\theobjnum
& I can describe the origin of the transmission line equation, understand the method of developing a solution, and interpret the solution.
\stepcounter{objnum}
\\\hline 

\thechapnum.\theobjnum
& I can explain the physical meaning of and calculate the following quantities: propagation constant, phase constant, attenuation constant, characteristic impedance, VSWR, reflection coefficient, and transmission coefficient.
\stepcounter{objnum}
\\\hline 

\thechapnum.\theobjnum
& I can use the transmission line equation to calculate the voltage/current at an arbitrary point on a specified transmission line.
\stepcounter{objnum}
\\\hline 

\thechapnum.\theobjnum
& I can use the transmission line equation to match a load to a source using a length of transmission line.
\stepcounter{objnum}
\\\hline 

\thechapnum.\theobjnum
& I can identify the various quantities represented on the Smith Chart.
\stepcounter{objnum}
\\\hline 

\thechapnum.\theobjnum
& I can use a Smith Chart to design a matching network for a given source and load, using series components, parallel components, shorted stubs, and open stubs.
\stepcounter{objnum}
\\\hline 

\thechapnum.\theobjnum
&  I can plot the transient response of a load connected to a source through a transmission line.
\stepcounter{objnum}
\\\hline 

\thechapnum.\theobjnum
& I can measure the reflection coefficient across a range of frequencies and interpret the measurement using a network analyzer.
\setcounter{objnum}{1}
\stepcounter{chapnum}
\\\hline 

%%%%%%%%%%%%%%%%%%%%%%%%%%%%%%%%%%%%%%%%%%%%%%%%%%%%%%%%%%%%%%%

\rowcolor{\chrowcolor}\multicolumn{2}{|c|}{Chapter {\thechapnum} - Vector Analysis}
\\\hline

\thechapnum.\theobjnum
& I can recognize and apply common vector calculus operations, such as: gradient, divergence, curl, and Laplacian in the three major coordinate systems.
\stepcounter{objnum}
\\\hline 

\thechapnum.\theobjnum
& I can calculate surface and volume integrals in rectangular, cylindrical, and spherical coordinate systems. 
\stepcounter{objnum}
\\\hline 

\thechapnum.\theobjnum
& I can differentiate between Stoke's Theorem and the Divergence Theorem and apply them to various problems to simplify integrals. \vspace{48pt}
\setcounter{objnum}{1}
\stepcounter{chapnum}
\\ \hline 

%%%%%%%%%%%%%%%%%%%%%%%%%%%%%%%%%%%%%%%%%%%%%%%%%%%%%%%%%%%%%%%

\rowcolor{\chrowcolor}\multicolumn{2}{|c|}{Chapter {\thechapnum} - Electrostatics}
\\\hline

\thechapnum.\theobjnum
&  I can apply Coloumb's Law to use direct integration to calculate the electric field due to charge density in varoius geometries, including rectangular, cylindrical, and spherical coordinate systems. 
\stepcounter{objnum}
\\\hline 

\thechapnum.\theobjnum
&  I can apply the integral and differential forms of Gauss' Law to calculate the electric field due to charge density in various geometries, including rectangular, cylindrical, and spherical coordinate systems. 
\stepcounter{objnum}
\\\hline 

\thechapnum.\theobjnum
&  I can  calculate the electric field from the scalar electric potential in a variety of geometries.
\stepcounter{objnum}
\\\hline 

\thechapnum.\theobjnum
&  I can apply the electric field boundary conditions to calculate electric fields on both sides of a material boundary. 
\setcounter{objnum}{1}
\stepcounter{chapnum}
\\ \hline 

%%%%%%%%%%%%%%%%%%%%%%%%%%%%%%%%%%%%%%%%%%%%%%%%%%%%%%%%%%%%%%%

\rowcolor{\chrowcolor}\multicolumn{2}{|c|}{Chapter {\thechapnum} - Steady Current and Conductivity}
\\\hline

\thechapnum.\theobjnum
&  I can calculate the total current through a surface area given a current distribution. 
\stepcounter{objnum}
\\ \hline 

\thechapnum.\theobjnum
&  I can use Ohm's Law for electromagnetics to calculate the relationship between a current distribution and an electric field. 
\stepcounter{objnum}
\\ \hline 

\thechapnum.\theobjnum
&  I can articulate and identify what makes a material a good conductor and how conductivity relates to resistance. 
\setcounter{objnum}{1}
\stepcounter{chapnum}
\\ \hline 

%%%%%%%%%%%%%%%%%%%%%%%%%%%%%%%%%%%%%%%%%%%%%%%%%%%%%%%%%%%%%%%

\rowcolor{\chrowcolor}\multicolumn{2}{|c|}{Chapter {\thechapnum} - Magnetostatics}
\\\hline

\thechapnum.\theobjnum
&  I can use the Biot-Savart law to simple magnetostatic systems to find the magnetic field at an observation point. 
\stepcounter{objnum}
\\ \hline 

\thechapnum.\theobjnum
&  I can articulate the meaning of both the integral form and point form of Gauss' Law for magnetism.
\stepcounter{objnum}
\\ \hline 

\thechapnum.\theobjnum
& I can use Ampere's Law to calculate the magnetic field at an observation point for simple magnetostatic systems.
\stepcounter{objnum}
\\ \hline 

\thechapnum.\theobjnum
&  I can apply the magnetic field boundary conditions to calculate magnetic fields on both sides of a material boundary.  
\setcounter{objnum}{1}
\stepcounter{chapnum}
\\ \hline 

%%%%%%%%%%%%%%%%%%%%%%%%%%%%%%%%%%%%%%%%%%%%%%%%%%%%%%%%%%%%%%%

\rowcolor{\chrowcolor}\multicolumn{2}{|c|}{Chapter {\thechapnum} - Dynamic Fields}
\\\hline

\thechapnum.\theobjnum
&  I can articulate the significance of the displacement current and calculate it, given a time-varying electric field.
\stepcounter{objnum}
\\ \hline 

\thechapnum.\theobjnum
& I can articulate the meaning of and calculate the relaxation time for a variety of materials. 
\stepcounter{objnum}
\\ \hline 

\thechapnum.\theobjnum
&  I know the forms of Maxwell's equations for time-varying fields and use them to calculate electric and magnetic fields for a variety of physical systems.  \vspace{48pt}
\setcounter{objnum}{1}
\stepcounter{chapnum}
\\ \hline 



%%%%%%%%%%%%%%%%%%%%%%%%%%%%%%%%%%%%%%%%%%%%%%%%%%%%%%%%%%%%%%%

\rowcolor{\chrowcolor}\multicolumn{2}{|c|}{Chapter {\thechapnum} - Plane Waves}
\\\hline

\thechapnum.\theobjnum
&  I understand the process of decoupling and finding solutions for Maxwell's equations for dynamic fields. 
\stepcounter{objnum}
\\ \hline 

\thechapnum.\theobjnum
& I can articulate what a Uniform Plane Wave (UPW) is and describe key properties that relate the electric and magnetic field. 
\stepcounter{objnum}
\\ \hline 

\thechapnum.\theobjnum
&  I can calculate the electric and magnetic fields for a UPW.
\stepcounter{objnum}
\\ \hline 

\thechapnum.\theobjnum
& I can articulate the concept of polarization and determine the polarization for various UPWs.
\stepcounter{objnum}
\\ \hline 

\thechapnum.\theobjnum
& I can determine the reflection and transmission of UPW at a boundary for Transverse Electric (TE) and Transverse Magnetic (TM) waves, using boundary conditions.
\setcounter{objnum}{1}
\stepcounter{chapnum}
\\ \hline 
\end{xltabular}			
\end{center}
\end{\myfontsize}
%-----------------------------------------------------------
